\section{Mathematical Model}

Based on the physical considerations discussed above, the coupled electromechanical behavior 
can be described at the continuum level by a system of partial differential equations.

The mechanical response of the material is governed by the equations of elasticity,

\begin{equation}
    \nabla \cdot \boldsymbol{\sigma}(u) = 0,
\end{equation}

where $u$ denotes the displacement field and $\boldsymbol{\sigma}(u)$ is the stress tensor, 
which depends on the strain.

The electrical response is described by the conductivity equation,

\begin{equation}
    \nabla \cdot \left( \sigma(\varepsilon(u)) \nabla V \right) = 0,
\end{equation}

where $V$ is the electrical potential and the conductivity $\sigma$ depends on the strain tensor 
$\varepsilon(u)$.

Together, these equations define a coupled system in which mechanical deformation influences 
electrical conduction through the strain-dependent conductivity. From a physical perspective, 
this coupling represents an effective description of how deformation alters the underlying 
electronic transport properties of the material.

This formulation provides the basis for computational modeling of the piezoresistive effect 
within a continuum framework.